This appendix gives the full definitions of the geometry metrics used in the
main text. The main paper groups these metrics by geometry channel: class
structure, radial norm geometry, covariance spectrum, local density, and
subspace structure. These metrics are not treated as competing explanations in
isolation; they are complementary fingerprints of the penultimate
representation.

\subsection{Class Structure and Neural Collapse Metrics}
\label{app:class-structure-nc}

Let \(\mu_c\) denote the class mean of ID training features and let
\[
    \bar{\mu}
    =
    \frac{1}{K}\sum_{c=1}^{K}\mu_c,
    \qquad
    \tilde{\mu}_c
    =
    \mu_c-\bar{\mu}.
\]
For class \(c\), let \(\Sigma_c\) be the class covariance and let
\[
    \Sigma_W
    =
    \frac{1}{K}\sum_{c=1}^{K}\Sigma_c
\]
be the pooled within-class covariance. The between-class covariance is
\[
    \Sigma_B
    =
    \frac{1}{K}\sum_{c=1}^{K}
    \tilde{\mu}_c \tilde{\mu}_c^\top .
\]
A standard NC1-style diagnostic is
\[
    \mathrm{NC1}
    =
    \frac{1}{K}\operatorname{tr}
    \left(\Sigma_W \Sigma_B^\dagger\right),
\]
where \(\dagger\) denotes the Moore--Penrose pseudoinverse. Smaller values
indicate stronger within-class collapse relative to between-class variation.

NC2 measures the angular structure of the centered class means
\(\{\tilde{\mu}_c\}_{c=1}^K\), typically by comparing their normalized Gram
matrix to the simplex ETF target. NC3 measures the alignment between classifier
weights \(w_c\) and centered class means \(\tilde{\mu}_c\). This is the class
structure channel most directly connected to CTM classifier-feature cosine
scores. We also report direct summaries such as WithinVar, the average
within-class dispersion, and InterDist, the average distance between class
means.

\subsection{Radial Norm Statistics}
\label{app:radial-norm-statistics}

We decompose each feature into norm and direction:
\[
    f_\theta(x)=r(x)u(x),
    \qquad
    r(x)=\|f_\theta(x)\|_2,
    \qquad
    \|u(x)\|_2=1.
\]
We measure the mean, standard deviation, quantiles, and class-wise dispersion of
\(r(x)\). These statistics test whether an optimizer changes feature magnitude
globally, by class, or by sample. Radial geometry can affect MaxLogit, Energy,
Mahalanobis, kNN, and GMM scores. Cosine-based scores such as CTM are
invariant to feature norm, and NeCo projection ratios normalize by the
feature norm.

\subsection{Covariance Spectrum and Conditioning}
\label{app:covariance-spectrum-conditioning}

For each class covariance, we consider the eigendecomposition
\[
    \Sigma_c
    =
    U_c
    \operatorname{diag}(\lambda_{c,1},\ldots,\lambda_{c,d})
    U_c^\top .
\]
We summarize the covariance spectrum by trace, leading eigenvalue, anisotropy,
effective rank, condition number, and cumulative explained variance. These
metrics test whether an optimizer spreads features unevenly across directions,
creates ill-conditioned covariance estimates, or changes the stability of
inverse-covariance detectors.

\subsection{Local Density Diagnostics}
\label{app:local-density-diagnostics}

kNN detectors do not assume a Gaussian covariance model. They read local
neighborhood spacing in the ID feature bank. Let
\[
    d_k(f_\theta(x),\mathcal{F}_{\mathrm{tr}})
\]
denote the distance from \(f_\theta(x)\) to its \(k\)-th nearest neighbor in the
training feature bank \(\mathcal{F}_{\mathrm{tr}}\). We compare the
distributions of this distance for ID and OOD samples. When feature labels are
available, we can additionally separate same-class and cross-class
nearest-neighbor distances for ID data. This reveals whether an optimizer
changes within-class local density or boundary-neighborhood geometry.

\subsection{Subspace Diagnostics}
\label{app:subspace-diagnostics}

Subspace-aware scores such as NeCo and PCA reconstruction error read whether a
feature lies in the principal ID subspace. Let \(P\in\mathbb{R}^{d\times r}\)
contain the top \(r\) principal directions of the ID training features. Then
\(PP^\top\) is the projection matrix onto the selected ID subspace. We measure
retained dimension, explained variance, projection norm, and residual norm:
\[
    \|P^\top f_\theta(x)\|_2,
    \qquad
    \|(I-PP^\top)f_\theta(x)\|_2.
\]
These quantities test whether an optimizer keeps ID features in a compact,
task-relevant subspace or spreads them into residual directions.
